\chapter{Arquitectura del Sistema}
\label{ch:arquitectura}

SIPREH sigue una arquitectura de tres capas débilmente acopladas: (i) almacenamiento en la nube, (ii) backend de procesamiento y API REST, y (iii) frontend web reactivo. Cada capa se comunica exclusivamente a través de la API HTTP versionada bajo \texttt{/api/v1}.

\section{Descripción y estructura del SIPREH}
\label{sec:arquitectura_detalle}

La Figura~\ref{fig:arquitectura_sipreh} muestra el flujo de datos desde las fuentes originales hasta el usuario final.

\begin{figure}[H]
\centering
\begin{tikzpicture}[
    every node/.style={font=\scriptsize\sffamily},
    src/.style={
        draw=violet, fill=violet!10, text=black,
        rounded corners=3pt, minimum width=2.6cm, minimum height=1.0cm,
        align=center, line width=0.8pt},
    cld/.style={
        draw=blue!60!black, fill=blue!8, text=black,
        rounded corners=3pt, minimum width=3.1cm, minimum height=1.0cm,
        align=center, line width=0.8pt},
    bck/.style={
        draw=teal, fill=teal!10, text=black,
        rounded corners=3pt, minimum width=3.1cm, minimum height=1.0cm,
        align=center, line width=0.8pt},
    frt/.style={
        draw=orange!80!black, fill=orange!8, text=black,
        rounded corners=3pt, minimum width=3.1cm, minimum height=1.0cm,
        align=center, line width=0.8pt},
    usr/.style={
        draw=green!60!black, fill=green!8, text=black,
        rounded corners=14pt, minimum width=14cm, minimum height=0.90cm,
        align=center, line width=0.8pt},
    band/.style={
        draw=gray!35, fill=gray!5, rounded corners=5pt,
        line width=0.7pt, inner sep=7pt},
    btitle/.style={
        fill=none, draw=none, font=\scriptsize\bfseries\sffamily,
        text=gray!55, anchor=north west},
    arr/.style={draw=gray!55, ->, line width=1.1pt},
]

%% ── ROW 1: FUENTES DE DATOS (y = 0) ────────────────────────────
\node[src] (era5l)  at ( 0.0, 0) {ERA5-Land\\0.1°\,·\,ECMWF};
\node[src] (imerg)  at ( 3.5, 0) {IMERG\\0.1°\,·\,NASA};
\node[src] (chirps) at ( 7.0, 0) {CHIRPS\\0.05°\,·\,UCSB};
\node[src] (ideam)  at (10.5, 0) {IDEAM\\29 estaciones};

\begin{scope}[on background layer]
  \node[band, fit=(era5l)(imerg)(chirps)(ideam)] (b1) {};
\end{scope}
\node[btitle] at (b1.north west) {Fuentes de Datos};

%% ── ROW 2: ALMACENAMIENTO R2 (y = -3.2) ────────────────────────
\node[cld] (hist)  at ( 1.75,-3.2) {Parquet Histórico\\SPI · SPEI · RAI\\EDDI · PDSI};
\node[cld] (pred)  at ( 5.50,-3.2) {Parquet Predicción\\\texttt{prediction\_main}\\12 horizontes};
\node[cld] (hydro) at (10.50,-3.2) {Parquet Hidrológico\\SDI · SRI · MFI\\DDI · HDI};

\begin{scope}[on background layer]
  \node[band, fit=(hist)(pred)(hydro)] (b2) {};
\end{scope}
\node[btitle] at (b2.north west) {Almacenamiento en la nube --- Cloudflare R2};

%% Flechas R1 → R2 (|- = vertical primero, luego horizontal)
\draw[arr] (era5l.south)  |- (hist.north);
\draw[arr] (imerg.south)  |- (hist.north);
\draw[arr] (chirps.south) |- (pred.north);
\draw[arr] (ideam.south)  -- (hydro.north);

%% ── ROW 3: BACKEND (y = -6.4) ──────────────────────────────────
\node[bck] (duckdb) at ( 1.75,-6.4) {Motor DuckDB\\Consultas 1D/2D\\Agg. cuencas};
\node[bck] (cache)  at ( 5.50,-6.4) {Caché Redis\\+ in-memory\\TTL 900 s};
\node[bck] (apibox) at (10.50,-6.4) {API REST\\FastAPI · /api/v1\\Auth JWT};

\begin{scope}[on background layer]
  \node[band, fit=(duckdb)(cache)(apibox)] (b3) {};
\end{scope}
\node[btitle] at (b3.north west) {Backend --- Python 3.11 · Uvicorn};

%% Flechas R2 → R3
\draw[arr] (hist.south)  -- (duckdb.north);
\draw[arr] (pred.south)  |- (duckdb.north);
\draw[arr] (hydro.south) -- (apibox.north);
\draw[arr] (duckdb.east) -- (cache.west);
\draw[arr] (cache.east)  -- (apibox.west);

%% ── ROW 4: FRONTEND (y = -9.6) ─────────────────────────────────
\node[frt] (mapa)   at ( 1.75,-9.6) {Mapa Interactivo\\Leaflet.js\\3 niveles zoom};
\node[frt] (charts) at ( 5.50,-9.6) {Gráficas\\uPlot (1D)\\Canvas (predicción)};
\node[frt] (admpan) at (10.50,-9.6) {Panel Admin\\Archivos · Sync\\Usuarios};

\begin{scope}[on background layer]
  \node[band, fit=(mapa)(charts)(admpan)] (b4) {};
\end{scope}
\node[btitle] at (b4.north west) {Frontend --- Next.js 14 · React 18};

%% Flechas R3 → R4 (apibox → tres componentes frontend)
%% Bus horizontal a y≈-7.15 (por debajo del band backend)
\draw[arr] (apibox.south) -- ++(0,-0.45) -| (mapa.north);
\draw[arr] (apibox.south) -- ++(0,-0.45) -| (charts.north);
\draw[arr] (apibox.south) -- (admpan.north);

%% ── ROW 5: USUARIO (y = -11.5) ─────────────────────────────────
\node[usr] (user) at (5.75,-11.5)
    {Usuario Final \quad
     (Investigador · Gestor de recursos hídricos · Administrador)};

\draw[arr] (b4.south) -- (user.north);

\end{tikzpicture}
\caption{Arquitectura en capas del SIPREH. Los datos fluyen verticalmente desde las
fuentes hasta el usuario final a través del almacenamiento en la nube, el backend y
el frontend.}
\label{fig:arquitectura_sipreh}
\end{figure}

El flujo de datos dentro del sistema sigue estos pasos:

\begin{enumerate}
    \item \textbf{Ingesta:} Los datos de ERA5-Land, IMERG y CHIRPS se procesan externamente por el pipeline y se almacenan como archivos Parquet en el bucket R2. Los datos del IDEAM se almacenan como Parquet hidrológico.

    \item \textbf{Descarga y caché de disco:} En la primera petición sobre un archivo, el backend lo descarga desde R2 y lo almacena localmente (hasta 10 GB). Las peticiones siguientes leen desde disco.

    \item \textbf{Procesamiento:} DuckDB ejecuta consultas SQL sobre los archivos Parquet. Soporta tres patrones: 1D (serie de tiempo por celda), 2D (campo espacial por fecha) y agregación por cuenca (promedio ponderado por área).

    \item \textbf{Caché de resultados:} Los resultados se almacenan en Redis con TTL de 900 s. Los accesos en caché devuelven respuestas en tiempos de un solo dígito de milisegundos ($\sim$50$\times$ más rápido).

    \item \textbf{API REST:} FastAPI expone los resultados bajo \texttt{/api/v1} con routers para autenticación, histórico, predicción, datos hidrológicos, administración y dashboard.

    \item \textbf{Visualización:} Next.js consume la API y renderiza el mapa (Leaflet.js), las gráficas 1D (uPlot) y las gráficas de predicción (Canvas 2D).
\end{enumerate}

\begin{definition}
\textbf{Agregación por cuenca:} el backend pondera el aporte de cada celda por la fracción de su área que intersecta el polígono de la cuenca. Las fracciones se precomputan al iniciar el sistema, por lo que la agregación en tiempo de consulta se reduce a un promedio ponderado.
\end{definition}
